\chapter{Cadre conceptuel et exigences de conception}

\section{Nature du problème traité}
La couche étudiée dans ce mémoire ne vise pas à produire directement des ordres
ou une stratégie complète, mais à transformer des données de marché en signaux
interprétables. Cette distinction est structurante. Un signal exprime une
information directionnelle ou conditionnelle sur l'état du marché, alors qu'une
stratégie ajoute des règles d'entrée et de sortie, des contraintes de risque,
des règles de taille de position et une logique d'exécution.

Sur le plan fonctionnel, cette séparation permet d'isoler le problème de la
\emph{qualité informative} du signal du problème plus large de la \emph{mise en
oeuvre d'une décision de trading}. Sur le plan méthodologique, elle évite de
masquer une faiblesse de signal derrière une ingénierie aval plus complexe.

\section{Architecture logique retenue}
Le périmètre retenu suit une chaîne ordonnée :
\[
\text{Données OHLCV validées} \rightarrow \text{Génération et validation des signaux}
\rightarrow \text{Analyse détaillée} \rightarrow \text{Restitution synthétique}
\]

Cette architecture impose trois niveaux de responsabilité.
\begin{enumerate}
\item Le premier niveau est celui de la \textbf{production méthodologique} : il
regroupe la génération des variantes, leur évaluation hors échantillon, leur
filtrage et leur agrégation.
\item Le deuxième niveau est celui de la \textbf{lecture analytique} : la page
Signaux permet de descendre du consensus global vers les familles puis vers les
représentants individuels.
\item Le troisième niveau est celui de la \textbf{lecture décisionnelle} : le
tableau de bord hiérarchise les titres et synthétise les résultats par horizon.
\end{enumerate}

Cette décomposition est cohérente avec les principes exposés par López de Prado,
qui insiste sur la nécessité de transformer les données brutes en variables
informatives avant toute phase de décision ou d'apprentissage \cite{deprado2018}.

\section{Principes méthodologiques directeurs}
Trois principes gouvernent l'ensemble de la couche signal.

\subsection{Primauté du hors échantillon}
Aucun résultat n'est retenu sur la seule base d'une performance observée sur la
période ayant servi au calcul préparatoire. Toute variante est jugée au travers
de fenêtres de test chronologiquement postérieures aux fenêtres de mise en
contexte. Cette exigence découle directement des recommandations de Pardo sur
l'analyse Walk-Forward et sur la nécessité de séparer nettement calibration et
évaluation \cite{pardo2008}.

\subsection{Robustesse avant performance apparente}
Le projet ne cherche pas le meilleur chiffre ponctuel, mais une structure de
signal capable de conserver un comportement acceptable à travers plusieurs
périodes. Autrement dit, le critère central n'est pas seulement le rendement
moyen, mais la combinaison de plusieurs dimensions : rendement ajusté du risque,
stabilité inter-périodes, cohérence et maîtrise des pertes. Cette logique répond
au problème du \emph{multiple testing}, souligné dans les travaux de Bailey et
López de Prado \cite{bailey2014}.

\subsection{Explicabilité des résultats}
Un score global n'a de valeur opérationnelle que s'il peut être relié à ses
composantes. L'architecture retenue impose donc qu'un signal agrégé soit
traçable jusqu'aux familles, puis jusqu'aux variantes représentatives. Cette
contrainte d'explicabilité détermine à la fois la structure du pipeline et celle
de l'interface utilisateur.

\section{Importance du calibrage par horizon}
Le traitement par horizon constitue l'un des choix les plus importants de la
couche signal. En effet, un paramètre pertinent pour une lecture de court terme
n'a pas la même signification économique pour une lecture de moyen ou de long
terme. Il n'est donc pas méthodologiquement correct d'évaluer une même grille de
paramètres sur tous les horizons.

Le projet distingue trois horizons, chacun associé à une thèse temporelle et à
une géométrie d'évaluation propres.

\begin{table}[H]
\centering
\caption{Paramètres d'évaluation par horizon}
\label{tab:horizon-eval}
\begin{tabular}{p{2.6cm}p{2.3cm}p{2.3cm}p{2.3cm}p{3cm}}
\toprule
\textbf{Horizon} & \textbf{Train} & \textbf{Test} & \textbf{Pas} & \textbf{Historique maximal} \\
\midrule
Court terme & 252 barres & 63 barres & 63 barres & 5 ans \\
Moyen terme & 504 barres & 126 barres & 126 barres & 10 ans \\
Long terme & 756 barres & 252 barres & 252 barres & 20 ans \\
\bottomrule
\end{tabular}
\end{table}

Le Tableau~\ref{tab:horizon-eval} montre que la notion d'horizon ne modifie pas
seulement l'affichage final ; elle transforme la profondeur historique utilisée,
la longueur des fenêtres de test et la fréquence de réévaluation. Ainsi, le
court terme cherche des signaux réactifs sur des structures rapides, tandis que
le long terme privilégie une lecture plus lente et plus robuste, compatible avec
une thèse d'investissement plus structurelle.

Le calibrage des familles suit la même logique. Pour les indicateurs fondés sur
un regard direct sur le passé, comme les moyennes mobiles, l'horizon s'exprime
par l'allongement de la fenêtre de calcul. Pour les oscillateurs bornés, tels
que le RSI, le changement d'horizon ne peut pas se réduire à une simple
augmentation du paramètre de période ; il faut également adapter les seuils de
déclenchement afin de conserver un pouvoir discriminant. Cette idée sera
reprise dans le chapitre méthodologique.

\section{Contraintes fonctionnelles du projet}
Le système étudié devait satisfaire plusieurs contraintes simultanées.
\begin{itemize}
\item Les données de marché devaient être ramenées à une forme OHLCV fiable et
reproductible.
\item Le coût des transactions devait être intégré dès l'évaluation du signal,
afin d'éviter des conclusions artificiellement optimistes.
\item La restitution devait rester lisible pour un utilisateur métier, sans
supprimer la possibilité d'audit analytique.
\item Les résultats devaient pouvoir être recalculés de manière stable pour un
symbole et un horizon donnés.
\end{itemize}

\section{Invariants conservés sur l'ensemble du périmètre}
À l'issue de cette analyse, quatre invariants peuvent être formulés.
\begin{enumerate}
\item Un score affiché ne doit jamais contourner l'évaluation hors échantillon.
\item Une famille ne peut influencer la synthèse qu'après passage par les étapes
de robustesse, de filtrage et de réduction de redondance.
\item Un signal global doit rester cohérent avec les contributions des familles
qui le composent.
\item L'utilisateur doit pouvoir passer d'une lecture synthétique à une lecture
détaillée sans rupture de sens.
\end{enumerate}

Ces invariants donnent au rapport sa cohérence interne : ils justifient la
structure du chapitre méthodologique, l'architecture de la page Signaux et la
forme de la restitution dans le tableau de bord.
